\chapter{The spectral compressor}\label{ch:compressor}

\begin{onething}
Before the transformer runs, a $1\times1$ convolution with BatchNorm replaces each pixel's 165
normalised bands with 32 learned mixtures, and afterwards a second $1\times1$ convolution, deliberately
without any normalisation, expands 32 back to 165; the pair is a linear bottleneck, so every
reconstructed spectrum lies in a 32-dimensional space.
\end{onething}

The spectral compressor is the only structural difference between Indradhanu and its thermal sibling
Chakshu. It is a short module, 59 lines in
\file{app/foundation_models/components/spectral_compressor.py}, most of them explanation. It is also
the module whose mathematics sets a limit on everything downstream, and the reason Part~IV of this book
(\cref{ch:ceiling}) exists. In this chapter you will compute what it does to one pixel by hand, learn
BatchNorm properly, count its parameters, and prove that the pair of convolutions can never reproduce
more than 32 independent spectral directions.

%======================================================================
\section{The problem: 165 channels into a 32-wide transformer}\label{sec:compressor:problem}

Indradhanu's transformer is a small SegFormer configuration (``B0''), whose first stage works with
token vectors of 32 numbers (\cref{ch:patchembed}). Hand it the 165 normalised bands directly and
three things go wrong.

\paragraph{The first layer becomes bloated.} Stage~1 begins with a $4\times4$ convolution from the input
channels to 32 (\cref{ch:shapes}). Its weights number $C_{\text{in}}\times32\times4\times4$:
\begin{align*}
\text{with } C_{\text{in}} = 165:&\quad 165\times32\times16 = 5{,}280\times16 = 84{,}480,\\
\text{with } C_{\text{in}} = 32:&\quad 32\times32\times16 = 1{,}024\times16 = 16{,}384.
\end{align*}
The ratio is $84{,}480/16{,}384 \approx 5.16$: five times the weights, on the layer that runs at the
finest resolution.

\paragraph{Most of the 165 numbers are near-duplicates.} Neighbouring bands on the common grid are
10~nm apart (\cref{ch:sensors}). Over 10~nm almost nothing changes in a natural spectrum, so bands~62
and~63 carry nearly the same value in nearly every pixel. You pay full price for information you
already have.

\paragraph{Every intermediate tensor is five times bigger.} A \shape{(B, 165, 128, 128)} activation holds
$165/32\approx5.2$ times as many numbers as \shape{(B, 32, 128, 128)}. Memory limits the batch size, and
the batch size limits training speed.

The obvious fix, ``pick 32 of the 165 bands'', throws away whatever the unpicked bands knew. The fix
the model uses is to \emph{mix}: build each of the 32 channels as a weighted blend of all 165 bands, with
the blend learned.

\begin{analogy}
A sound engineer mixes a 48-track recording down to a stereo pair. They do not keep two tracks and
throw away 46; each output channel is a blend of every track, set so that what matters survives. A
listener's speakers later turn the stereo pair back into a room full of sound. The compressor is the
mixdown, the decompressor is the playback, and here the mixing-desk settings are learned. The analogy
also tells you what is lost: two instruments panned identically can never be separated again from the
stereo pair. \Cref{sec:compressor:rank} makes that precise.
\end{analogy}

\coursefig[\linewidth]{iw-10-compressor}{The spectral bottleneck around the SegFormer. A $1\times1$
convolution with BatchNorm mixes each pixel's 165 normalised bands into 32 channels (5{,}376
parameters); the spatial model runs on those 32; a second $1\times1$ convolution, without
normalisation, expands them back to 165 (5{,}445 parameters). (Course figure IW~10.1.)}{fig:compressor:overview}

%======================================================================
\section{The compressor is a matrix applied at every pixel}\label{sec:compressor:matrix}

The code is short enough to quote whole:
\begin{lstlisting}
class SpectralCompressor(nn.Module):
    def __init__(self, in_channels, compressed_channels):
        super().__init__()
        self.compress = nn.Conv2d(in_channels, compressed_channels, kernel_size=1)
        self.norm = nn.BatchNorm2d(compressed_channels)

    def forward(self, x):          # x: (B, 165, H, W)
        x = self.compress(x)       # x: (B, 32, H, W)
        x = self.norm(x)           # x: (B, 32, H, W)
        return x

class SpectralDecompressor(nn.Module):
    def __init__(self, compressed_channels, out_channels):
        super().__init__()
        self.decompress = nn.Conv2d(compressed_channels, out_channels, kernel_size=1)

    def forward(self, x):          # x: (B, 32, H, W)
        return self.decompress(x)  # -> (B, 165, H, W)
\end{lstlisting}

\Cref{ch:shapes} showed that a $1\times1$ convolution is a linear layer applied independently at every
pixel (\cref{eq:shapes:onebyone}). So the convolution part of the compressor is, for every pixel
$(i,j)$,
\begin{equation}\label{eq:compressor:conv}
\vect{z}_{:,i,j} = \mat{W}\,\vect{x}_{:,i,j} + \vect{b},\qquad
\mat{W}\in\R^{32\times165},\ \vect{b}\in\R^{32},
\end{equation}
where $\vect{x}_{:,i,j}$ is the pixel's 165 normalised band values (\cref{ch:bandpipeline}). Row $d$ of
$\mat{W}$ is the recipe for compressed channel $d$: how much of each band goes into it. The same
$\mat{W}$ is used at every pixel of every patch; the compressor never looks at a neighbour.

\begin{worked}[label={we:compressor:mix}]{Four bands into two}
Pretend there are 4 bands and 2 compressed channels. One pixel, already normalised:
\[
\vect{x} = (1.2,\ 1.0,\ -0.4,\ -0.6)\T .
\]
Suppose training arrived at these recipes, with zero biases:
\[
\mat{W} = \begin{pmatrix}0.5 & 0.5 & 0.0 & 0.0\\ 0.0 & 0.0 & 0.5 & 0.5\end{pmatrix}.
\]
Channel~0 uses row~0:
\begin{align*}
z_0 &= 0.5\times1.2 + 0.5\times1.0 + 0.0\times(-0.4) + 0.0\times(-0.6)\\
&= 0.60 + 0.50 + 0 + 0 = 1.10.
\end{align*}
Channel~1 uses row~1:
\begin{align*}
z_1 &= 0.0\times1.2 + 0.0\times1.0 + 0.5\times(-0.4) + 0.5\times(-0.6)\\
&= 0 + 0 - 0.20 - 0.30 = -0.50.
\end{align*}
So $\vect{z} = (1.10,\ -0.50)\T$. Four numbers became two, and the two still describe the shape of the
spectrum: bright at the short-wavelength end, dark at the long. The real compressor does the same with
165 inputs and 32 outputs, with recipes found by gradient descent (\cref{ch:trainloop}).
\end{worked}

\Cref{fig:shapes:conv1x1} in \cref{ch:shapes} draws this example.

%======================================================================
\section{BatchNorm}\label{sec:compressor:bn}

The compressor's second line is \code{nn.BatchNorm2d(32)}. Batch normalisation \citep{ioffe2015batch}
appears throughout the convolutional parts of this repository, and it behaves differently in training
and in use, which has caused real trouble in this project (\cref{ch:finetune}). It deserves a careful
treatment.

\subsection{What it computes in training}

\begin{definition}[Batch normalisation, training mode]
Let $z_{b,d,i,j}$ be channel $d$ of the input to \code{BatchNorm2d}, for patch $b$ of the batch at pixel
$(i,j)$. For each channel $d$ separately, \emph{batch normalisation}\index{batch normalisation}
\begin{enumerate}
\item computes the mean and variance of that channel over the whole batch and every pixel, that is over
  $M = B\cdot H\cdot W$ values:
  \begin{equation}\label{eq:compressor:bnstats}
  \mu_d = \frac{1}{M}\sum_{b,i,j} z_{b,d,i,j},\qquad
  \sigma_d^2 = \frac{1}{M}\sum_{b,i,j} \bigl(z_{b,d,i,j} - \mu_d\bigr)^2;
  \end{equation}
\item standardises every value of the channel with them, adding a small $\varepsilon$ (here $10^{-5}$)
  so that a constant channel does not divide by zero:
  \begin{equation}\label{eq:compressor:bnstd}
  \hat{z}_{b,d,i,j} = \frac{z_{b,d,i,j} - \mu_d}{\sqrt{\sigma_d^2 + \varepsilon}};
  \end{equation}
\item rescales and shifts with two learned parameters per channel, the \emph{scale}
  $\gamma_d$\index{batch normalisation!scale and shift} and the \emph{shift} $\beta_d$:
  \begin{equation}\label{eq:compressor:bnaffine}
  u_{b,d,i,j} = \gamma_d\,\hat{z}_{b,d,i,j} + \beta_d .
  \end{equation}
\end{enumerate}
\end{definition}

After step~2 each channel has mean $0$ and variance (almost exactly) $1$ over the batch. Step~3 lets
training choose a different mean ($\beta_d$) and spread ($\gamma_d$) if that suits the next layer
better; they start at $\gamma_d = 1$ and $\beta_d = 0$. The variance in \cref{eq:compressor:bnstats}
divides by $M$, not $M-1$: it is the \emph{biased} variance.

\begin{worked}[label={we:compressor:bn}]{BatchNorm on four values}
One channel; the batch and pixels together give four values $z = 1,\ 2,\ 3,\ 6$. Take $\gamma = 1$,
$\beta = 0$. Work on paper.

Mean:
\begin{align*}
1 + 2 + 3 + 6 &= 12,\\
\mu &= 12/4 = 3.
\end{align*}
Deviations from the mean: $1-3 = -2$, $2-3 = -1$, $3-3 = 0$, $6-3 = 3$.
Squares: $4,\ 1,\ 0,\ 9$. Their sum:
\[
4 + 1 + 0 + 9 = 14.
\]
Biased variance:
\[
\sigma^2 = 14/4 = 3.5.
\]
Denominator:
\[
\sqrt{3.5 + 0.00001} = \sqrt{3.50001} = 1.87083.
\]
Standardised values, one per deviation:
\begin{align*}
-2 / 1.87083 &= -1.0690,\\
-1 / 1.87083 &= -0.5345,\\
0 / 1.87083 &= 0,\\
3 / 1.87083 &= 1.6036.
\end{align*}
With $\gamma=1$, $\beta=0$ these are the outputs. With $\gamma = 2$, $\beta = 0.5$ instead, each is
doubled and shifted: $2\times(-1.0690) + 0.5 = -1.6381$, $2\times(-0.5345)+0.5 = -0.5690$,
$2\times0+0.5 = 0.5$, $2\times1.6036 + 0.5 = 3.7071$. PyTorch gives the same eight numbers\tagC.
\end{worked}

\subsection{Running statistics and evaluation mode}

At inference there is often one patch, or a batch whose make-up is accidental; normalising each scene by
its own batch statistics would make a pixel's output depend on which other patches happened to be
processed with it. So during training BatchNorm also keeps \emph{running statistics}\index{running
statistics}, exponential moving averages of the batch statistics, stored as buffers (not parameters;
they are not trained by gradient descent):
\begin{align}
\text{running\_mean}_d &\leftarrow (1-m)\,\text{running\_mean}_d + m\,\mu_d,\label{eq:compressor:rm}\\
\text{running\_var}_d &\leftarrow (1-m)\,\text{running\_var}_d + m\,s_d^2,\label{eq:compressor:rv}
\end{align}
with momentum $m = 0.1$ (PyTorch's default, used here) and starting values $0$ and $1$. Note a
PyTorch detail: the running variance is fed the \emph{unbiased} variance
$s_d^2 = \frac{1}{M-1}\sum(z-\mu_d)^2$, while the normalisation in training uses the biased one.

In \emph{evaluation mode} (\code{model.eval()}), BatchNorm stops computing batch statistics and uses the
running ones:
\begin{equation}\label{eq:compressor:bneval}
u_{b,d,i,j} = \gamma_d\,\frac{z_{b,d,i,j} - \text{running\_mean}_d}{\sqrt{\text{running\_var}_d + \varepsilon}} + \beta_d .
\end{equation}
Each output now depends only on its own input. Training mode and evaluation mode give the same answer
only when the running statistics match the statistics of the data you feed.

\begin{worked}[label={we:compressor:bneval}]{One training step of running statistics, then evaluation}
Continue \cref{we:compressor:bn} with a fresh layer (running mean $0$, running variance $1$). The batch
mean was $\mu = 3$. The unbiased variance divides the same sum of squares by $M-1 = 3$:
\[
s^2 = 14/3 = 4.66667.
\]
Update with $m = 0.1$:
\begin{align*}
\text{running\_mean} &= 0.9\times0 + 0.1\times3 = 0 + 0.3 = 0.3,\\
\text{running\_var} &= 0.9\times1 + 0.1\times4.66667 = 0.9 + 0.46667 = 1.36667.
\end{align*}
Now switch to evaluation mode and feed the same four values ($\gamma=1$, $\beta=0$). The denominator is
$\sqrt{1.36667 + 0.00001} = 1.16905$, and
\begin{align*}
(1 - 0.3)/1.16905 &= 0.7/1.16905 = 0.5988,\\
(2 - 0.3)/1.16905 &= 1.7/1.16905 = 1.4542,\\
(3 - 0.3)/1.16905 &= 2.7/1.16905 = 2.3096,\\
(6 - 0.3)/1.16905 &= 5.7/1.16905 = 4.8758.
\end{align*}
In training mode the value $6$ became $1.6036$; in evaluation mode, after one step, it becomes $4.8758$.
After thousands of steps on similar data the running statistics converge and the two modes agree. When
the data in use differ from the training data, they do not. PyTorch reproduces all of these
numbers\tagC.
\end{worked}

\subsection{The buffers travel inside the checkpoint}

\code{BatchNorm2d(32)} has $2\times32 = 64$ trainable parameters ($\gamma$ and $\beta$) and three
buffers: \code{running_mean} (32 values), \code{running_var} (32) and \code{num_batches_tracked} (one
counter)\tagC. That is $65$ values that are saved in the checkpoint but are not parameters. Together
with the $660$ values of normalisation statistics they account for the gap between Indradhanu's
5{,}506{,}629 parameters and the 5{,}507{,}354 values in its state dictionary (\cref{ch:recon}).

\subsection{The convolution's bias is redundant in front of BatchNorm}

Here is a small fact worth deriving, because it surprises people. Add a constant $b_d$ to every value of
channel $d$ before BatchNorm. In training mode,
\begin{align*}
\text{mean of } (z + b_d) &= \mu_d + b_d
  && \text{a constant shifts the mean by itself}\\
(z + b_d) - (\mu_d + b_d) &= z - \mu_d,
  && \text{the shift cancels in the deviation}
\end{align*}
and the variance, built from deviations, is unchanged too. So the standardised value
\cref{eq:compressor:bnstd} does not depend on $b_d$ at all. The 32 biases of the compressor's
convolution have no effect on its training-mode output, receive zero gradient, and stay at their initial
values; the running mean absorbs them, so they cancel in evaluation mode as well. Adding 10 to every bias
leaves the training-mode output unchanged, as the compute script confirms\tagC. They are
harmless, and they are still counted as parameters.

\subsection{Why the compressor has a BatchNorm}

The compressor's job is to hand the encoder a stable, well-scaled input. Whatever mixtures $\mat{W}$
learns, their scale can drift during training; BatchNorm resets every compressed channel to a controlled
mean and spread, so the first patch embedding always sees inputs of a predictable size. The module's own
docstring: ``The compressor's job is to hand the encoder a stable distribution''.

%======================================================================
\section{The decompressor, and why it must have no normalisation}\label{sec:compressor:decomp}

The decompressor maps the decoder's 32-channel output back to 165 bands at every pixel:
\begin{equation}\label{eq:compressor:decomp}
\hat{\vect{x}}_{:,i,j} = \mat{V}\,\hat{\vect{z}}_{:,i,j} + \vect{c},\qquad
\mat{V}\in\R^{165\times32},\ \vect{c}\in\R^{165}.
\end{equation}
It is \emph{not} the inverse of the compressor. $\mat{W}$ is $32\times165$ and has no inverse; $\mat{V}$
is a separate set of learned weights, trained together with everything else so that the whole chain
reconstructs well. Nothing ties $\mat{V}$ to $\mat{W}$.

It has no BatchNorm and no activation function, and the docstring warns against ``tidying'' that:
``the decompressor has NO norm and NO activation \ldots the decompressor's output feeds denormalisation
and must stay free to represent any value''. Here is what each would break.

\paragraph{A BatchNorm would erase what the batch has in common.} Training-mode BatchNorm forces every
output band to mean $\beta_d$ over the batch. But the batch's true mean in a band is information: a batch
of patches over open water has normalised near-infrared values well below zero in every pixel. Take two
water pixels whose true normalised value in some band is $-1.5$, and suppose the decoder predicts
exactly that, $-1.5$ and $-1.5$. BatchNorm computes mean $-1.5$, deviations $0$ and $0$, and outputs
$\beta$ for both. The reconstruction now says ``average brightness'' where the truth was ``dark water'',
and the residual is large at every pixel of the batch, ordinary or not. The answer would also depend on
which other patches shared the batch.

\paragraph{An activation would clip.} A ReLU sets every negative value to zero. In normalised space a
negative value means ``below this band's average''; after a ReLU no band could ever be reconstructed
darker than its average.

The same reasoning applies to the final layer of the old convolutional decoder, whose docstring
(\file{app/foundation_models/components/spatial_decoder.py}) explains that BatchNorm there ``would force
the output to zero-mean unit-variance per channel - but real thermal data is NOT zero-mean or
unit-variance''.

\begin{keypoint}
The asymmetry is the design: the compressor normalises because the encoder needs a stable input; the
decompressor does not, because its output must be free to take any value.
\end{keypoint}

%======================================================================
\section{A linear bottleneck: everything passes through 32 directions}\label{sec:compressor:rank}

Now the mathematics that matters most. Whatever happens between them, the compressor and decompressor
impose a hard limit on what the model can output.

\subsection{Rank, briefly}

The \emph{rank}\index{rank (of a matrix)} of a matrix is the number of linearly independent columns it
has, equivalently the dimension of the set of all vectors $\mat{A}\vect{u}$ it can produce
(\cref{ch:linalg}). A $165\times32$ matrix has only 32 columns, so its rank is at most 32.

\subsection{The output lies in a 32-dimensional space}

Look again at \cref{eq:compressor:decomp}. Whatever the encoder and decoder do, and they are nonlinear
and spatial, their output at pixel $(i,j)$ is some vector $\hat{\vect{z}}_{:,i,j}\in\R^{32}$. Write its
entries $\hat{z}_1,\dots,\hat{z}_{32}$ and the columns of $\mat{V}$ as $\vect{v}_1,\dots,\vect{v}_{32}$.
Then
\begin{align}
\hat{\vect{x}}_{:,i,j} &= \mat{V}\hat{\vect{z}}_{:,i,j} + \vect{c}
  && \text{\cref{eq:compressor:decomp}}\notag\\
&= \hat{z}_1\vect{v}_1 + \hat{z}_2\vect{v}_2 + \dots + \hat{z}_{32}\vect{v}_{32} + \vect{c}.
  && \text{a matrix times a vector is a mix of its columns}\label{eq:compressor:span}
\end{align}
Every reconstructed spectrum, at every pixel of every scene, is the fixed vector $\vect{c}$ plus a
mixture of the same 32 fixed spectra $\vect{v}_1,\dots,\vect{v}_{32}$. In the language of
\cref{ch:linalg}: it lies in a 32-dimensional affine subspace of $\R^{165}$. Any part of a real spectrum
that points outside that subspace can never be reconstructed, however good the transformer is.

There is a matching limit on the way in. $\mat{W}$ has 32 rows, so at most 32 independent combinations
of the bands reach the transformer; any change to a pixel's spectrum along a direction $\vect{n}$ with
$\mat{W}\vect{n} = \vect{0}$ is invisible to everything after the compressor.

\subsection{The chain without the transformer}

If the transformer were removed and BatchNorm ignored, the chain would be
\begin{align*}
\hat{\vect{x}} &= \mat{V}(\mat{W}\vect{x} + \vect{b}) + \vect{c}
  && \text{compress, then decompress}\\
&= \mat{V}\mat{W}\,\vect{x} + (\mat{V}\vect{b} + \vect{c}).
  && \text{multiply out the bracket}
\end{align*}
The $165\times165$ matrix $\mat{V}\mat{W}$ has rank at most 32. The best such matrix for reconstructing
data, in the least-squares sense, projects onto the data's 32 leading principal components
(\cref{ch:linalg}); the trained compressor is not constrained to find that, because it is trained jointly
with a network that also uses spatial context.

\begin{worked}[label={we:compressor:lost}]{A 3-to-2-to-3 bottleneck and what it loses}
Three bands, two compressed channels, zero biases:
\[
\mat{W} = \begin{pmatrix}1 & 0 & 0\\ 0 & 0.5 & 0.5\end{pmatrix},\qquad
\mat{V} = \begin{pmatrix}1 & 0\\ 0 & 1\\ 0 & 1\end{pmatrix}.
\]
Pixel $\vect{x} = (0.2,\ 0.4,\ 0.6)\T$. Compress:
\begin{align*}
z_1 &= 1\times0.2 + 0\times0.4 + 0\times0.6 = 0.2,\\
z_2 &= 0\times0.2 + 0.5\times0.4 + 0.5\times0.6 = 0 + 0.2 + 0.3 = 0.5.
\end{align*}
Decompress, row by row of $\mat{V}$:
\begin{align*}
\hat{x}_1 &= 1\times0.2 + 0\times0.5 = 0.2,\\
\hat{x}_2 &= 0\times0.2 + 1\times0.5 = 0.5,\\
\hat{x}_3 &= 0\times0.2 + 1\times0.5 = 0.5.
\end{align*}
Residual:
\[
\vect{x} - \hat{\vect{x}} = (0.2-0.2,\ 0.4-0.5,\ 0.6-0.5) = (0,\ -0.1,\ 0.1).
\]
The product is
\[
\mat{V}\mat{W} = \begin{pmatrix}1&0&0\\0&0.5&0.5\\0&0.5&0.5\end{pmatrix},
\]
whose last two rows are equal, so its rank is $2$\tagC. What was lost? The difference between bands~2
and~3. Check that the direction $\vect{n} = (0,\ 1,\ -1)\T$ is invisible to the compressor:
\begin{align*}
\mat{W}\vect{n} &= \bigl(1\times0 + 0\times1 + 0\times(-1),\ \ 0\times0 + 0.5\times1 + 0.5\times(-1)\bigr)\T\\
&= (0,\ 0)\T .
\end{align*}
The pixel's residual $(0, -0.1, 0.1)$ is exactly $-0.1\,\vect{n}$. Every pixel whose bands~2 and~3
differ leaves that difference in its residual, ordinary or anomalous. If an anomaly's signature were a
band-2-versus-band-3 contrast, this model would see it in the residual, but so would it see the
ordinary variation along $\vect{n}$, at every pixel.
\end{worked}

\Cref{we:compressor:lost} is the whole story in miniature. The bottleneck is lossy; the bet is that
what it loses is mostly redundancy and noise. Training decides which 32 directions to keep, and keeps
the ones that matter for reconstructing ordinary terrain. A narrow absorption feature that ordinary
terrain rarely shows may lie largely outside them. \Cref{ch:ceiling} measures that cost and is the
reason the project later built a cross-attention stem in place of the compressor.

%======================================================================
\section{Counting the parameters}\label{sec:compressor:count}

Use $\text{in}\times\text{out}+\text{out}$ for a $1\times1$ convolution (\cref{ch:shapes}) and
$2\times\text{channels}$ for a BatchNorm.

\begin{worked}[label={we:compressor:count}]{The spectral bottleneck's 10,821 parameters}
Compressor convolution, $165\to32$:
\begin{align*}
165\times32 &= 5{,}280,\\
5{,}280 + 32 &= 5{,}312.
\end{align*}
Compressor BatchNorm:
\[
2\times32 = 64.
\]
Compressor total:
\[
5{,}312 + 64 = 5{,}376.
\]
Decompressor convolution, $32\to165$:
\begin{align*}
32\times165 &= 5{,}280,\\
5{,}280 + 165 &= 5{,}445.
\end{align*}
Both:
\[
5{,}376 + 5{,}445 = 10{,}821.
\]
The \code{torchinfo} table in \file{research/model_break_down/05_hyperspectral_segformer_mae.md} lists
\code{Conv2d: 2-1} with 5{,}312, \code{BatchNorm2d: 2-2} with 64 and \code{Conv2d: 2-19} with
5{,}445\tagM{research/model_break_down/05_hyperspectral_segformer_mae.md}, and summing
\code{parameters()} of the real modules gives 5{,}376 and 5{,}445\tagC. Out of 5{,}506{,}629
parameters, the bottleneck is
\[
10{,}821 / 5{,}506{,}629 = 0.00197,
\]
under $0.2\%$ of the model.
\end{worked}

It is almost free, and it saves the $84{,}480 - 16{,}384 = 68{,}096$ weights of
\cref{sec:compressor:problem} in stage~1 alone, plus a factor of five in memory at full resolution.

%======================================================================
\section{It is MNF-shaped, but it is not MNF}\label{sec:compressor:mnf}

\Cref{ch:localrx} met the Minimum Noise Fraction transform, MNF \citep{green1988transformation}: a
classical linear projection from many bands to few, found by an eigen-decomposition that ranks
directions by signal-to-noise ratio. Structurally it is the same kind of object as $\mat{W}$: a
matrix of shape (few, many). The differences are in how the matrix is found and what it serves.

\begin{center}
\begin{tabular}{lll}
\toprule
& MNF & \texttt{SpectralCompressor}\\
\midrule
matrix found by & eigen-decomposition & gradient descent\\
optimises & signal-to-noise ratio & reconstruction loss\\
runs & offline, per scene & learned once, all scenes\\
learned with the network & no & yes\\
\bottomrule
\end{tabular}
\end{center}

The last row is the real advantage: MNF optimises a proxy, while the compressor optimises the thing the
model is graded on, and adapts while the rest of the network adapts to it. The module docstring calls
it ``a trainable analogue of MNF''.

%======================================================================
\section{Why $D$ is 32, not 8 and not 128}\label{sec:compressor:width}

\textbf{Not much larger:} stage~1 of the encoder embeds tokens at 32 numbers, so feeding it far more
than 32 channels inflates the first convolution only to compress again immediately.
\textbf{Not much smaller:} squeeze too hard and real spectral structure is destroyed before the
transformer sees it; by \cref{eq:compressor:span} the model then cannot get spectral shape right however
long it trains.

The project's training configuration (the file
\file{configs/hyperspectral_segformer_exp_1.json}) records this as a diagnosis rule:
\begin{lstlisting}
"SAM_stuck_above_0.3_after_epoch_30":
    "Compressor bottleneck too narrow (try compressed_channels=32)"
\end{lstlisting}
SAM is the spectral-angle error
(\cref{ch:losses}). That is exactly the change made between versions.

\begin{drift}
\begin{tabular}{lll}
& v0.1.0 & v0.2.0\\
compressed channels $D$ & 24$^{a}$ & 32$^{b}$\\
final validation loss & 0.07694 & 0.04349\\
\end{tabular}

$^{a}$\tagM{configs/hyperspectral_segformer_exp_1.json}\\
$^{b}$\tagM{configs/hyperspectral_segformer_exp_2.json}

Validation losses from course part IW10 (quoting a checkpoint manifest that is not in the repository
tree) and, for v0.2.0, \file{research/model_break_down/05_hyperspectral_segformer_mae.md}. Total
parameters: 5.21~M at $D=24$, 5.51~M at $D=32$\tagM{research/model_break_down/05_hyperspectral_segformer_mae.md}.
The shipped model since 2026-10 is Indradhanu~V2 (\cref{ch:shipping}).
\end{drift}

\begin{sourcenote}[: the code default is 24, the shipped model is 32]
Both \code{HyperspectralSegFormerMAE.__init__}
(\file{app/foundation_models/components/hyperspectral_seg_former_mae.py}, line~55) and
\code{HyperspectralSegFormerMAEConfig} (\file{app/models/training/training_config.py}, line~89) default
\code{compressed_channels} to 24, and the class docstring's example says ``(e.g. 24)''. The v0.2.0
model and its successors are built with 32 from the configuration file. If you construct the class
without arguments you get the old, narrower bottleneck, with
$(165\times24 + 24 + 48) + (24\times165 + 165) = 4{,}032 + 4{,}125 = 8{,}157$ compressor and
decompressor parameters instead of 10{,}821\tagC.
\end{sourcenote}

%======================================================================
\begin{confusions}
\confusion{Is the decompressor the inverse of the compressor?}{No. A $32\times165$ matrix has no
inverse. The decompressor is an independent set of learned weights; nothing ties it to the compressor
except that both are trained to make the chain reconstruct well.}
\confusion{Does the compressor look at neighbouring pixels?}{No. Its kernel is $1\times1$, so every
pixel is processed alone. All spatial reasoning happens in the transformer.}
\confusion{Is the compression lossy?}{Yes. By \cref{eq:compressor:span} every output spectrum lies in a
32-dimensional space. The bet is that what is lost is mostly redundancy and noise; validation loss and
detection scores test the bet.}
\confusion{Why is it called spectral compression when the picture stays $128\times128$?}{Because only the
band axis is compressed. Height and width are untouched.}
\confusion{Does BatchNorm normalise each pixel's spectrum?}{No. It normalises each \emph{channel} across
the batch and all pixels. Normalising each token across its own features is LayerNorm
(\cref{ch:patchembed}).}
\confusion{Are BatchNorm's running mean and variance parameters?}{No. They are buffers, updated by
\cref{eq:compressor:rm,eq:compressor:rv}, not by gradients. They are saved in the checkpoint and used in
evaluation mode.}
\end{confusions}

\begin{checkyourself}
\item\label{ex:compressor:1} Compute the parameter count of \code{Conv2d(165, 24, kernel_size=1)} and of the
whole v0.1.0 compressor (with its \code{BatchNorm2d(24)}).
\item\label{ex:compressor:2} A BatchNorm channel receives the batch values $2,\ 4,\ 4,\ 6$ with
$\gamma = 1$, $\beta = 0$, $\varepsilon$ negligible. Compute the training-mode outputs. Then, starting
from running mean $0$ and running variance $1$ with momentum $0.1$, compute the running statistics after
this one step.
\item\label{ex:compressor:3} Give the two reasons the decompressor has no BatchNorm and no activation,
each with a one-line example.
\item\label{ex:compressor:4} In \cref{we:compressor:lost}, find the reconstruction of the pixel
$\vect{x} = (0.3,\ 0.5,\ 0.5)\T$ and its residual. Explain the result using $\vect{n}$.
\item\label{ex:compressor:5} Show that adding a constant $b$ to every input value of a BatchNorm channel
leaves its training-mode output unchanged, using the values of \cref{we:compressor:bn} with $b = 10$.
\item\label{ex:compressor:6} How is the compressor like MNF, and how is it different? What symptom
suggests $D$ is too small, and what was done about it?
\end{checkyourself}

\nextchapter{ch:patchembed}
